% YY14200B
% Einschätzung der Indikation zur Wirbelsäulenimmobilisation
% 14.0
% 2013-11-30
:ref:`m-einschaetzung-indikation-ws-immobilisation`
\label{proc:YY14200B}
Eine Immobilisation der Wirbelsäule ist notwendig bei
Traumapatienten, bei denen folgendes zutrifft:

\begin{enumerate}
  -   **Typischer Unfallmechanismus**
  \begin{compactitem}
    -   Hochgeschwindigkeitsunfall
    -   Sturz $\geq$ 3\Meter* oder 3fache Patientengröße
    -   Penetrierende Verletzungen im Bereich der
    Wirbelsäule
    -   Sportverletzungen im Kopf- oder Nackenbereich
    -   Trauma nach Sprung ins Wasser
    -   Stauchungsverletzung der HWS (z. B. Schlag auf
    Kopf)
    -   Suizidversuch durch Erhängen
    -   Bewusstloser Traumapatient
  \end{compactitem}
  -   **Unklarer Unfallmechanismus**, bei dem einer der
  folgenden Punkte zutrifft:
  \begin{compactitem}
    -   Spinaler Schmerz oder Druckschmerz über der
    Wirbelsäule
    -   Auffälliger Befund bei Untersuchung von Motorik
    und Sensibilität\footnote{sofern nicht anders plausibel
    erklärbar}
    -   Patientenangaben kann nicht vertraut
    werden
    \begin{compactitem}
      -   Stresssituation
      -   Schädel- oder zerebrale Verletzung
      -   Auffälliger Bewusstseinszustand
      -   Intoxikation
      -   Ablenkende Verletzungen
    \end{compactitem}
  \end{compactitem}
  -   **Verdacht**
\end{enumerate}

\noindent Solange es noch unklar ist, ob eine WS-Immobilisation
notwendig ist, muss eine manuelle Fixierung der HWS
durchgeführt werden!

:cite:`ITLSde6,ScholzNotfallmedizin2,SiewertChirurgie8`\ \Commentary[Konsensus]{Review 2011-01-11}